% =============================================================================
%  Article de recherche - Simulateur d'entraînement au contrôle aérien
%  augmenté par IA : validation systématique et étude comparative multi-modèles
%  Auteur : Nicolas Marano - URCA / supercalculateur ROMEO (NVIDIA GH200)
%
%  Compilation :  docs/article/build.sh  (copie les figures, génère les
%  tableaux depuis bench/results/*.json puis pdflatex x2 via WSL).
%  AUCUN chiffre n'est saisi à la main : tables/ et numbers.tex sont générés
%  par gen_tables.py depuis les JSON de résultats.
% =============================================================================
\documentclass[10pt,a4paper]{article}
\usepackage[T1]{fontenc}
\usepackage[utf8]{inputenc}
\usepackage[french]{babel}
\usepackage[margin=2.2cm]{geometry}
\usepackage{graphicx}
\usepackage{booktabs}
\usepackage{amsmath}
\usepackage{tikz}
\usetikzlibrary{positioning,arrows.meta,shapes.geometric,fit}
\usepackage{xcolor}
\usepackage{caption}
\captionsetup{font=small,labelfont=bf}
\usepackage[colorlinks=true,linkcolor=blue!50!black,citecolor=blue!50!black,
            urlcolor=blue!50!black,pdftitle={Simulateur ATC augmenté par IA},
            pdfauthor={Nicolas Marano}]{hyperref}

% GENERE par gen_tables.py — ne pas editer a la main
\newcommand{\nCPADesaccords}{0}
\newcommand{\nCPAErrMax}{0,05}
\newcommand{\nCPAGeoms}{1 000 000}
\newcommand{\nConflitDcpaMax}{0,082}
\newcommand{\nConflitDcpaMoy}{0,018}
\newcommand{\nConflitGarantie}{100,0\,\%}
\newcommand{\nConflitTerrMax}{0,7}
\newcommand{\nEELatMoy}{2,6\,s}
\newcommand{\nEELatPceQuinze}{3,2\,s}
\newcommand{\nEETexte}{88\,\%}
\newcommand{\nEEVoix}{76\,\%}
\newcommand{\nEEn}{25}
\newcommand{\nLLMBestAcc}{81,9\,\%}
\newcommand{\nLLMBestName}{Mistral-7B}
\newcommand{\nLLMSweetAcc}{81,9\,\%}
\newcommand{\nLLMSweetLat}{0,8\,s}
\newcommand{\nLLMSweetName}{Mistral-7B}
\newcommand{\nLLMWorstAcc}{3,4\,\%}
\newcommand{\nLLMWorstName}{Llama-3.2-1b}
\newcommand{\nParserHorsGram}{70,8\,\%}
\newcommand{\nParserInGrammar}{100,0\,\%}
\newcommand{\nScalingCinq}{105}
\newcommand{\nScalingDeuxCents}{67}
\newcommand{\nScenBest}{71,8\,\%}
\newcommand{\nScenParser}{100,0\,\%}
\newcommand{\nTTSKokoroRTF}{0,336}
\newcommand{\nTTSKokoroWERVHF}{21,4\,\%}
\newcommand{\nWERGainRel}{60\,\%}
\newcommand{\nWERLoraAtcoVhf}{28,5\,\%}
\newcommand{\nWERLoraUwbVhf}{19,2\,\%}
\newcommand{\nWERVanillaAtcoVhf}{71,9\,\%}
                        % macros numériques GÉNÉRÉES

\newcommand{\syst}[1]{\texttt{#1}}
\newcommand{\ic}[2]{[#1 ; #2]}
% inclusion robuste : cadre placeholder si la figure n'est pas encore generee
\newcommand{\figpng}[2]{\IfFileExists{figures/#1.png}%
  {\includegraphics[width=#2\linewidth]{figures/#1}}%
  {\fbox{\parbox[c][3cm][c]{.8\linewidth}{\centering\small figure \texttt{#1}\\(bench non exécuté)}}}}

\title{\bfseries Simulateur d'entraînement au contrôle aérien augmenté par IA :\\
validation systématique et étude comparative multi-modèles\\
de la chaîne vocale STT\,$\rightarrow$\,LLM\,$\rightarrow$\,TTS}
\author{Nicolas Marano\\
\small Université de Reims Champagne-Ardenne, Centre de calcul ROMEO (NVIDIA GH200)\\
\small \texttt{maranoniko@gmail.com}}
\date{Juillet 2026}

\begin{document}
\maketitle

% =============================================================================
\begin{abstract}
\noindent
Nous présentons un simulateur d'entraînement au contrôle aérien dont la boucle
radio est entièrement pilotée par la voix : l'élève contrôleur parle sur un
canal VHF simulé, la parole est transcrite (Whisper adapté au domaine par
LoRA), interprétée en ordres structurés par un grand modèle de langage
contraint par une base de connaissances OACI, validée par des gardes
déterministes (bornes physiques, graphe secteur), exécutée dans le simulateur
de trafic BlueSky, puis collationnée par une voix de synthèse dégradée canal
radio. Toutes les briques d'inférence respectent le contrat d'API
OpenAI-compatible : le même système fonctionne indifféremment sur
supercalculateur (ROMEO, GH200), sur GPU grand public local (RTX 4070) ou sur
API cloud, par simple changement de configuration.

Nous validons chaque étage par des mesures contrôlées et reproductibles :
(i)~géométrie de conflit : \nCPAGeoms{} géométries Monte-Carlo sur 5 graines,
\nCPADesaccords{} désaccord de décision avec une vérité numérique
indépendante, et confrontation à BlueSky sur 200 rencontres
($F_1=1{,}0$) ; (ii)~reconnaissance vocale : WER sur corpus VHF réels
(ATCO2, UWB-ATCC), l'adaptation LoRA réduisant le WER ATCO2 de
\nWERVanillaAtcoVhf{} à \nWERLoraAtcoVhf{} ; (iii)~interprétation : comparaison appariée (McNemar exact, correction de Holm, rapports de cotes) de
cinq LLM locaux quantifiés et d'un parseur à règles sur 116 clairances
annotées, dont des strates hors grammaire et adversariales ; le meilleur
modèle local (\nLLMBestName) atteint \nLLMBestAcc{} d'exactitude TrafScript
exacte, contre \nParserInGrammar{} pour le parseur à règles sur la
phraséologie standard mais \nParserHorsGram{} hors grammaire ;
(iv)~synthèse : coût temps réel et intelligibilité aller-retour ;
(v)~boucle complète : \nEEVoix{} de clairances parlées correctement
exécutées (témoin texte : \nEETexte), latence vocale totale moyenne de
\nEELatMoy{}. Deux campagnes de généralisation complètent l'étude : la même
batterie rejouée \emph{à protocole identique} contre la façade
supercalculateur (GH200, modèles pleine précision, tunnel SSH de
production), isolant l'effet de la quantification 4~bits et établissant la
parité STT entre points de déploiement ; et une validation avec locuteur
humain réel (\nHumMicro{} de réussite au micro brut contre \nEEVoix{} en
voix synthétique), dont l'analyse des échecs fonde une taxonomie
actionnable. La campagne historique de validation se ré-exécute à
l'identique (mêmes métriques champ à champ). L'ensemble du banc de mesure,
des corpus annotés et des scripts d'analyse est versionné avec le projet.
\end{abstract}

% =============================================================================
\section{Introduction}

La formation initiale des contrôleurs aériens repose sur des simulateurs dont
le coût d'exploitation est dominé par un facteur humain : chaque session
requiert des \emph{pseudo-pilotes}, opérateurs qui prononcent les
collationnements et exécutent les instructions dans le simulateur. Automatiser ce rôle (comprendre une clairance parlée, l'exécuter fidèlement, la collationner d'une voix réaliste) est un cas d'usage exigeant pour les
modèles de fondation : la phraséologie OACI est un langage contraint, le canal
VHF (300--3\,400~Hz) dégrade fortement le signal, et toute erreur
d'interprétation exécutée dans le simulateur fausse l'exercice pédagogique.

Ce travail décrit un système complet et son évaluation : une application web
d'entraînement (radar, strips, poste instructeur, exercices notés) adossée au
simulateur de trafic open source BlueSky~\cite{bluesky}, dont la boucle radio
est assurée par trois services d'inférence interchangeables au contrat
OpenAI-compatible. Nos contributions :

\begin{enumerate}
  \item une \textbf{architecture de sûreté en couches} : aucune sortie de LLM
        n'atteint le simulateur sans passer par une validation déterministe
        (bornes physiques, actions autorisées, waypoints du secteur), et toute
        panne de fournisseur est remontée visiblement (jamais de repli
        silencieux) ;
  \item un \textbf{banc de mesure de bout en bout}, versionné et seedé,
        couvrant les cinq étages (géométrie/simulateur, STT, LLM, TTS, boucle
        vocale complète) avec intervalles de confiance bootstrap, IC de Wilson
        et tests appariés de McNemar ;
  \item une \textbf{étude comparative multi-modèles reproductible sur GPU
        grand public} : cinq LLM open-weights quantifiés (1~à 7~milliards de
        paramètres) évalués \emph{à travers le code de production}, comparés
        à un parseur à règles de référence, sur un corpus annoté stratifié
        (phraséologie standard, paraphrases hors grammaire, bruit STT,
        français, cas mixtes valide$+$invalide, négatifs de sécurité) ;
  \item la \textbf{quantification de l'adaptation au domaine} : gain du LoRA
        ATC sur Whisper-small mesuré sur enregistrements VHF réels, effet du
        filtre passe-bande d'inférence, et coût (RTF) des points de
        déploiement local, cluster et cloud ;
  \item une \textbf{comparaison cluster/local à protocole strictement
        identique} (mêmes corpus, mêmes graines, même chaîne cliente) entre la
        façade GH200 (modèles pleine précision) et un GPU portable grand
        public (modèles quantifiés), ainsi qu'une \textbf{validation avec
        locuteur humain réel} de la boucle vocale complète.
\end{enumerate}

\paragraph{Questions de recherche.} L'évaluation est organisée autour de
cinq questions falsifiables :
\begin{description}
  \item[RQ1] Le prédicteur de conflit embarqué est-il \emph{exact} au sens
    d'une vérité numérique indépendante, et la garantie « conflit par
    construction » des exercices tient-elle statistiquement ?
  \item[RQ2] Quel gain de reconnaissance vocale l'adaptation LoRA
    apporte-t-elle sur des enregistrements VHF \emph{réels}, et le filtre
    passe-bande d'inférence du système est-il justifié ?
  \item[RQ3] Un LLM local remplace-t-il avantageusement un parseur à règles
    pour l'extraction d'ordres --- globalement, et hors de la grammaire du
    parseur ? Les différences entre modèles survivent-elles à une correction
    de comparaisons multiples ?
  \item[RQ4] Les gardes déterministes suffisent-elles à empêcher toute
    exécution d'ordre dangereux, même avec des modèles faibles ?
  \item[RQ5] La boucle vocale complète est-elle exploitable en conditions
    réalistes (canal radio dégradé, locuteur humain, fournisseur distant),
    à quelle latence, et où échoue-t-elle ?
\end{description}

% =============================================================================
\section{Système}
\label{sec:systeme}

\subsection{Architecture}

La figure~\ref{fig:archi} présente l'architecture. Le navigateur (React,
canvas radar 60~fps, push-to-talk) communique avec un serveur FastAPI par REST
et WebSocket ($\sim$8~Hz). Le simulateur BlueSky, non thread-safe, vit dans un
thread unique alimenté par une file de commandes. Les trois briques
d'inférence sont des services REST au contrat OpenAI standard
(\syst{/v1/audio/transcriptions}, \syst{/v1/chat/completions},
\syst{/v1/audio/speech}), configurés par variables d'environnement : changer de fournisseur, qu'il s'agisse de la façade auto-hébergée sur ROMEO, de la façade locale llama.cpp / faster-whisper / Kokoro (fournie avec le banc de mesure) ou d'une API cloud, ne modifie ni le code ni le comportement de validation.

\begin{figure}[t]
\centering
\begin{tikzpicture}[
  font=\scriptsize,
  box/.style={draw, rounded corners=2pt, align=center, inner sep=4pt,
              minimum height=7mm, fill=blue!6},
  ai/.style={box, fill=orange!12},
  guard/.style={box, fill=green!12},
  sim/.style={box, fill=purple!8},
  fleche/.style={-{Stealth[length=2mm]}, thick},
  node distance=4mm and 7mm]

\node[box] (ui) {Navigateur\\radar canvas · strips · PTT (V)};
\node[box, right=13mm of ui] (app) {FastAPI \syst{atc\_app}\\REST + WebSocket 8 Hz};
\node[ai, above right=5mm and 11mm of app] (stt) {STT\\\syst{/v1/audio/transcriptions}};
\node[ai, right=11mm of app] (llm) {LLM + KB OACI\\\syst{/v1/chat/completions}};
\node[ai, below right=5mm and 11mm of app] (tts) {TTS + VHF client\\\syst{/v1/audio/speech}};
\node[guard, below=6mm of app] (guard) {Validation déterministe\\bornes HDG/ALT/SPD · graphe secteur\\cohérence climb/descend};
\node[sim, below=6mm of guard] (sim) {BlueSky (thread unique)\\conflits CD natif + repli CPA\\vent · turbulence · zones};
\node[box, right=10mm of sim] (ex) {Exercices notés\\conflits garantis · score /100};

\draw[fleche] (ui) -- node[above]{audio WAV 16 kHz} (app);
\draw[fleche] (app) -- (stt);
\draw[fleche] (app) -- (llm);
\draw[fleche] (app) -- (tts);
\draw[fleche] (app) -- node[right]{ordres JSON}(guard);
\draw[fleche] (guard) -- node[right]{TrafScript} (sim);
\draw[fleche] (sim) -| (ex);
\draw[fleche] (app.170) -- node[above]{état + événements WS} (ui.10);

\node[draw, dashed, rounded corners, fit=(stt)(llm)(tts),
      label={[font=\scriptsize]90:{fournisseurs interchangeables (.env) : ROMEO GH200 · local RTX 4070 · cloud}}] {};
\end{tikzpicture}
\caption{Architecture. Les briques IA (orange) sont des services
OpenAI-compatibles interchangeables ; aucune sortie LLM n'atteint le
simulateur sans la couche de validation déterministe (vert).}
\label{fig:archi}
\end{figure}

\subsection{Sûreté en couches}
\label{sec:surete}

Le contrat d'interprétation impose un tableau JSON d'ordres
$\{$\syst{callsign}, \syst{action}$\in\{$HDG, ALT, SPD, ADDWPT$\}$,
\syst{value}$\}$. En aval du modèle, une bibliothèque pure et testée
(\syst{03\_bluesky\_connector}) applique : coercition stricte de type (rejet
des booléens, NaN, infinis), bornes physiques (HDG $[0;360]^\circ$, ALT
$[0;45\,000]$~ft, SPD $[0;350]$~kt), et validation des waypoints contre le
graphe secteur. L'application ajoute deux gardes : un indicatif absent du
radar produit un silence radio (réalisme assumé), et une incohérence
sémantique montée/descente (« climb » vers un niveau inférieur au niveau
courant) est rejetée avec message. Enfin, aucune panne de fournisseur n'est
rattrapée silencieusement : erreur typée $\rightarrow$ HTTP~502 $\rightarrow$
événement visible dans le journal de bord de l'interface.

\subsection{Adaptation au domaine}

La brique STT est \syst{openai/whisper-small} affiné par LoRA ($r=32$,
$\alpha=64$, cibles $q/v$-proj) sur les corpus radiotéléphoniques réels
UWB-ATCC et ATCOSIM~\cite{uwbatcc,atcosim}, avec un front-end reproduisant le
canal : passe-bande Butterworth d'ordre~6, 300--3\,400~Hz, et bruit additif
(SNR 5--20~dB) à l'entraînement. Le prompt LLM inline une base de connaissances OACI originale de 9 fiches ($\sim$2,5~Ko), couvrant les actions, les contextes informatifs à ignorer (QNH, transferts de fréquence) et la règle « en cas de doute, ne rien produire », ainsi que des indices NER (indicatif détecté).
Le collationnement TTS applique la phraséologie OACI (chiffres parlés,
alphabet radio, indicatifs en téléphonie) avec une voix stable par aéronef
(hachage CRC32 de l'indicatif dans un pool configurable) et la dégradation
VHF appliquée \emph{côté client}, donc identique quel que soit le fournisseur.

% =============================================================================
\section{Méthodologie}
\label{sec:methodo}

\paragraph{Principe.} Chaque étage est mesuré \emph{à travers le code de
production} (mêmes prompts, mêmes clients HTTP, même post-traitement), jamais
par une reconstitution. Tous les protocoles sont seedés et versionnés
(\syst{bench/}) ; chaque script écrit un JSON auto-porteur (protocole +
résultats) et les figures et tableaux de cet article en sont générés
automatiquement.

\paragraph{Statistiques.} Moyennes : IC à 95\,\% par bootstrap percentile
($B=10^4$, graine fixée)~\cite{efron}. Proportions : IC de
Wilson~\cite{wilson}. Comparaisons de systèmes sur corpus apparié : test
exact de McNemar~\cite{mcnemar} (binomial, bilatéral) pour les issues
binaires, avec correction de Holm-Bonferroni~\cite{holm} sur la famille
complète des paires et rapport de cotes $n_{01}/n_{10}$ comme taille
d'effet ; bootstrap apparié pour les différences de WER. Percentiles de
latence par interpolation linéaire. Aucune valeur de l'article n'est saisie
à la main : tableaux et macros numériques sont générés depuis les JSON de
mesure (\syst{docs/article/gen\_tables.py}), et l'environnement logiciel
exact est capturé au moment des campagnes (tableau~\ref{tab:versions}).

\begin{table}[t]\centering\small
\caption{Environnement logiciel exact des campagnes (reproductibilité).}
\label{tab:versions}
\begin{tabular}{ll}
\toprule
Composant & Version \\
\midrule
Python (bench) & 3.12.10 \\
OS & Windows 11 \\
torch & 2.6.0+cu124 \\
transformers & 5.12.1 \\
peft & 0.19.1 \\
datasets & 5.0.0 \\
faster\_whisper & 1.2.1 \\
ctranslate2 & 4.8.0 \\
llama\_cpp & 0.3.32 \\
kokoro\_onnx & ? \\
onnxruntime & 1.27.0 \\
jiwer & ? \\
numpy & 2.5.0 \\
scipy & 1.18.0 \\
soundfile & 0.14.0 \\
fastapi & 0.139.0 \\
requests & 2.34.2 \\
GPU (local) & NVIDIA GeForce RTX 4070 Laptop GPU \\
CUDA (torch) & 12.4 \\
Pilote NVIDIA & 591.74 \\
Python (app) & 3.12.10 \\
bluesky-simulator & 1.1.1 \\
numpy (app) & 2.4.6 \\
Modeles GGUF (Q4\_K\_M) & Llama-3.2-1B-Instruct-Q4\_K\_M (0.81 Go), Mistral-7B-Instruct-v0.3-Q4\_K\_M (4.37 Go), Qwen2.5-3B-Instruct-bartowski-Q4\_K\_M (1.93 Go), qwen2.5-1.5b-instruct-q4\_k\_m (1.12 Go), qwen2.5-3b-instruct-q4\_k\_m (2.10 Go) \\
\bottomrule
\end{tabular}
\end{table}


\paragraph{Corpus d'interprétation.} 116 clairances annotées en TrafScript
attendu \emph{après} validation : 68 cas historiques dans la grammaire du
parseur de référence (10 catégories dont 10 négatifs) et 48 cas étendus couvrant paraphrases hors grammaire, bruit type STT (hésitations, répétitions),
français parlé, nombres composés, immatriculations épelées, multi-ordres à
connecteurs, cas \emph{mixtes} (ordre valide $+$ ordre invalide dans la même
phrase, l'invalide devant être filtré et le valide conservé) et 14 négatifs
de sécurité supplémentaires (hors bornes, hors schéma, instructions
relatives, information pure). Un cas est réussi si le multi-ensemble
TrafScript produit est exactement celui attendu.

\paragraph{Corpus audio.} STT : 150 extraits par corpus, tirés au sort
(graine 42, durées 0,4--30~s) de deux corpus VHF \emph{réels}, ATCO2-1h
(871 extraits disponibles, domaine cible jamais vu à l'entraînement du LoRA)
et UWB-ATCC test (2\,822 extraits, domaine d'entraînement) ; les IC
bootstrap par extrait reflètent cette taille d'échantillon. Normalisation texte et WER
identiques au protocole historique du projet (comparabilité avec les chiffres
publiés). Boucle E2E : énoncés contrôleur synthétisés (voix Windows SAPI,
jamais utilisée ailleurs dans le banc), passés dans le canal radio simulé
(bruit SNR 12~dB puis passe-bande VHF).

\paragraph{Matériel et points de déploiement.} Étude locale : portable
Windows~11, RTX~4070 Laptop 8~Go, i9-13900H, 32~Go (versions logicielles
exactes au tableau~\ref{tab:versions}). LLM locaux : llama.cpp (CUDA),
GGUF Q4\_K\_M, température 0, graine 42 --- Llama-3.2-1B-Instruct,
Qwen2.5-1.5B-Instruct, Qwen2.5-3B-Instruct (deux conversions GGUF
indépendantes pour contre-vérification) et Mistral-7B-Instruct-v0.3.
Campagne cluster (août 2026) : la façade de production sur nœud GH200
(Whisper-small$+$LoRA, Mistral-7B-Instruct-v0.3 \emph{bf16}, XTTS voix
clonées), atteinte par le tunnel SSH de production, soumise aux mêmes
protocoles seedés --- Mistral-7B étant évalué des deux côtés, la paire
bf16/Q4\_K\_M isole l'effet de quantification à modèle constant. Les
mesures historiques de \syst{validation/results\_perf.json} (2026-07)
restent citées comme référence de première campagne.

% =============================================================================
\section{Résultats}
\label{sec:resultats}

\subsection{Géométrie de conflit et simulateur}
\label{sec:res-geo}

\paragraph{Prédiction CPA vs vérité numérique.} Sur \nCPAGeoms{} géométries
aléatoires (5 graines $\times$ 200\,000 ; positions $\pm$80~NM, caps
uniformes, vitesses sol 150--550~kt), la décision du prédicteur de conflit
embarqué (extrapolation linéaire, horizon 120~s, seuil 5~NM) coïncide avec
une recherche numérique indépendante sur grille fine ($\Delta t=0{,}05$~s,
horizon 300~s) dans \emph{tous} les cas : \nCPADesaccords{} désaccord sur
$10^6$ décisions. L'écart maximal sur la distance au point de rapprochement
est de \nCPAErrMax~NM, borné par la quantification d'affichage du champ
(arrondi au dixième de NM) ; la campagne historique, qui comparait les valeurs non arrondies, mesurait un écart maximal de
$5{,}52\times10^{-4}$~NM (figure~\ref{fig:cpa}).

\paragraph{Confrontation à BlueSky.} La campagne historique (200 rencontres
stabilisées, 55\,\% biaisées conflit, seuils 5~NM/120~s) donne précision,
rappel et $F_1$ de $1{,}0$ (113 conflits), MAE de $0{,}067$~NM sur $d_{CPA}$
et $0{,}71$~s sur $t_{CPA}$. Nous l'avons ré-exécutée intégralement : toutes
les métriques sont identiques champ à champ à celles versionnées (seules les
durées d'exécution diffèrent), ce qui établit la reproductibilité du
protocole.

\paragraph{Garantie de conflit des exercices.} Le générateur d'exercices
construit des paires d'aéronefs convergeant vers un point commun avec temps
d'arrivée égaux. Sur 2\,000 tirages, la garantie est vérifiée à
\nConflitGarantie{} : $d_{CPA}$ analytique maximal de \nConflitDcpaMax~NM
(moyenne \nConflitDcpaMoy~NM, très inférieure à la norme de séparation de
5~NM) et écart maximal entre l'instant de croisement annoncé et le CPA
calculé de \nConflitTerrMax~s (figure~\ref{fig:conflit}).

\paragraph{Montée en charge.} BlueSky headless avance 10~s de simulation avec
un facteur temps réel moyen de $\times$\nScalingCinq{} à 5 aéronefs et encore
$\times$\nScalingDeuxCents{} à 200 aéronefs (5 répétitions par point,
figure~\ref{fig:scaling}) : la simulation n'est jamais le facteur limitant de
la boucle vocale.

\begin{figure}[t]
\centering
\figpng{fig_cpa_multiseed}{.72}
\caption{Erreur de $d_{CPA}$ du prédicteur embarqué vs grille numérique fine,
par graine (200\,000 géométries chacune). L'erreur est bornée par l'arrondi
d'affichage (0,05~NM) ; la décision de prédiction est exacte
(0 désaccord / $10^6$).}
\label{fig:cpa}
\end{figure}

\begin{figure}[t]
\centering
\figpng{fig_conflit_garanti}{.62}
\caption{Garantie « conflit par construction » du générateur d'exercices :
distribution du $d_{CPA}$ analytique sur 2\,000 tirages seedés.}
\label{fig:conflit}
\end{figure}

\begin{figure}[t]
\centering
\figpng{fig_sim_scaling}{.7}
\caption{Montée en charge BlueSky headless (5 répétitions par point, IC
bootstrap à 95\,\%).}
\label{fig:scaling}
\end{figure}

\subsection{Reconnaissance vocale (STT)}
\label{sec:res-stt}

Le tableau~\ref{tab:stt} et la figure~\ref{fig:stt} comparent cinq systèmes
sur les deux corpus VHF réels, avec et sans passe-bande VHF d'inférence
(protocole de production : \emph{avec}). Trois enseignements :

\begin{enumerate}
\item \textbf{L'adaptation LoRA est le facteur dominant.} Sur ATCO2 (domaine
      jamais vu à l'entraînement), le LoRA réduit le WER de
      \nWERVanillaAtcoVhf{} (vanilla) à \nWERLoraAtcoVhf{}, gain relatif
      d'environ \nWERGainRel{} confirmé par bootstrap apparié (différence
      significative, l'IC de la différence excluant zéro). Sur UWB-ATCC
      (domaine d'entraînement), le WER tombe à \nWERLoraUwbVhf{}.
\item \textbf{Le passe-bande d'inférence est cohérent avec l'entraînement.}
      Il améliore ou laisse inchangé le modèle adapté (entraîné dans ce
      domaine), tandis que son effet sur les modèles génériques est faible :
      le choix de production (\syst{bandpass=True}) est validé.
\item \textbf{Le déploiement local est réaliste.} En greedy sur RTX 4070,
      tous les systèmes tournent bien en deçà du temps réel
      (figure~\ref{fig:sttrtf}) ; faster-whisper offre le meilleur débit
      brut, la voie transformers+LoRA offrant elle la meilleure exactitude.
\end{enumerate}

\begin{table}[t]\centering\small
\caption{STT : WER corpus (\%) et RTF par système, corpus VHF réels, sans/avec passe-bande VHF d'inférence. Normalisation du protocole historique du projet ; décodage greedy.}
\label{tab:stt}
\begin{tabular}{lrrrrrr}
\toprule
 & \multicolumn{3}{c}{atco2 (n=150)} & \multicolumn{3}{c}{uwb\_atcc (n=150)} \\
 & brut & VHF & RTF & brut & VHF & RTF \\
\midrule
hf-small & 71,1 & 71,9 & 0,018 & 101,2 & 91,9 & 0,023 \\
hf-small-lora & 29,7 & 28,5 & 0,017 & 19,3 & 19,2 & 0,020 \\
fw-tiny & 100,9 & 106,0 & 0,106 & 120,1 & 123,3 & 0,112 \\
fw-base & 88,1 & 90,1 & 0,091 & 107,4 & 110,8 & 0,148 \\
fw-small & 70,6 & 71,5 & 0,092 & 93,2 & 92,3 & 0,130 \\
\bottomrule
\end{tabular}
\end{table}


\begin{figure}[t]
\centering
\figpng{fig_stt_wer}{1}
\caption{WER par système sur corpus VHF réels (IC bootstrap 95\,\% du WER
moyen par extrait ; losanges : WER corpus). Protocole de normalisation
identique aux mesures historiques du projet.}
\label{fig:stt}
\end{figure}

\begin{figure}[t]
\centering
\figpng{fig_stt_rtf}{.75}
\caption{Facteur temps réel STT (log) : tous les systèmes sont largement
sous le temps réel sur GPU grand public.}
\label{fig:sttrtf}
\end{figure}

\subsection{Interprétation des clairances (LLM)}
\label{sec:res-llm}

Le tableau~\ref{tab:llm} et les figures~\ref{fig:llmacc}--\ref{fig:llmlat}
présentent l'étude comparative sur les 116 clairances. Lecture :

\begin{enumerate}
\item \textbf{Le parseur à règles est parfait dans sa grammaire, nul en
      dehors.} \nParserInGrammar{} sur la phraséologie standard,
      \nParserHorsGram{} sur les strates hors grammaire : c'est précisément
      l'écart que le LLM doit combler.
\item \textbf{L'échelle compte, mais pas de façon monotone.} L'exactitude
      globale croît de \nLLMWorstAcc{} (\nLLMWorstName) à \nLLMBestAcc{}
      (\nLLMBestName) ; les différences entre modèles sont testées par
      McNemar apparié (tableau~\ref{tab:mcnemar}). Une anomalie notable :
      Qwen2.5-3B s'effondre ($\approx$24\,\%) là où son cadet 1.5B excelle : il honore le message système (vérifié) mais viole le schéma
      (champ \syst{callsign} omis, clés JSON non citées). Nous avons exclu
      un artefact de conversion en ré-évaluant un second GGUF indépendant
      (bartowski) : résultat identique (23,3\,\%). Autre artefact
      \emph{réellement} rencontré puis corrigé : le gabarit de conversation
      Mistral-v0.3 de llama.cpp \emph{perd le message système} (vérifié par
      test dédié), faisant chuter Mistral-7B à 21\,\% ; la fusion
      système$\rightarrow$utilisateur du gabarit officiel restaure
      \nLLMBestAcc{}. Ces deux cas illustrent pourquoi l'évaluation doit se
      faire \emph{à travers la pile de déploiement réelle}.
\item \textbf{La sécurité tient d'abord aux gardes déterministes.} Les taux
      de rejet des négatifs par système (tableau~\ref{tab:llm}) montrent que
      même les petits modèles, qui hallucinent des ordres, sont largement
      rattrapés par la validation de bornes et de waypoints, illustration chiffrée de l'architecture en couches (\S\ref{sec:surete}).
\item \textbf{Latence exploitable en local.} Le meilleur compromis
      qualité/latence locale est \nLLMSweetName{} (\nLLMSweetAcc,
      \nLLMSweetLat{} par clairance en moyenne) ; Mistral-7B, identique au
      modèle ROMEO, sert de point de comparaison cluster/local.
\end{enumerate}

\begin{table}[t]\centering\small
\caption{Interprétation : exactitude TrafScript exacte (IC Wilson 95\,\%), par strate, rejet des négatifs de sécurité, JSON strict, latence et débit (116 clairances ; chaîne de production complète).}
\label{tab:llm}
\begin{tabular}{lrrrrrrr}
\toprule
Système & Globale & IC 95\,\% & Std. & Hors gr. & Nég. rejetés & JSON & Lat. moy. \\
\midrule
Parseur à règles & 87,9 & \ic{81}{93} & 100 & 71 & 96 & 100 & 0,00\,s \\
\midrule
Llama-3.2-1b & 3,4 & \ic{1}{9} & 1 & 6 & 17 & 22 & 2,81\,s \\
Mistral-7B & 81,9 & \ic{74}{88} & 85 & 77 & 75 & 100 & 0,82\,s \\
Qwen2.5-3b-bartowski & 23,3 & \ic{17}{32} & 21 & 27 & 96 & 69 & 0,42\,s \\
Qwen2.5-1.5b & 75,0 & \ic{66}{82} & 79 & 69 & 75 & 97 & 0,32\,s \\
Qwen2.5-3b & 24,1 & \ic{17}{33} & 19 & 31 & 100 & 46 & 0,34\,s \\
Mistral-7B ROMEO (bf16) & 83,6 & \ic{76}{89} & 84 & 83 & 75 & 98 & 1,95\,s \\
\bottomrule
\end{tabular}
\end{table}

\begin{table}[t]\centering\small
\caption{Tests exacts de McNemar entre systèmes appariés sur les 116 clairances ($n_{01}$ : A correct / B faux ; RC = rapport de cotes $n_{01}/n_{10}$ ; $p_{\mathrm{Holm}}$ : correction de Holm-Bonferroni sur la famille complète de 21 comparaisons ; * : $p_{\mathrm{Holm}}<0{,}05$).}
\label{tab:mcnemar}
\begin{tabular}{lrrrrr}
\toprule
Paire (A vs B) & $n_{01}$ & $n_{10}$ & RC & $p$ & $p_{\mathrm{Holm}}$ \\
\midrule
Parseur à règles vs Llama-3.2-1b & 98 & 0 & $\infty$ & $<10^{-4}$ & $<10^{-4}$* \\
Parseur à règles vs Mistral-7B & 19 & 12 & 1,6 & 0,2810 & 1,0000\phantom{*} \\
Parseur à règles vs Qwen2.5-3b-bartowski & 76 & 1 & 76,0 & $<10^{-4}$ & $<10^{-4}$* \\
Parseur à règles vs Qwen2.5-1.5b & 25 & 10 & 2,5 & 0,0167 & 0,1167\phantom{*} \\
Parseur à règles vs Qwen2.5-3b & 75 & 1 & 75,0 & $<10^{-4}$ & $<10^{-4}$* \\
Parseur à règles vs Mistral-7B ROMEO (bf16) & 18 & 13 & 1,4 & 0,4731 & 1,0000\phantom{*} \\
Llama-3.2-1b vs Mistral-7B & 0 & 91 & 0,0 & $<10^{-4}$ & $<10^{-4}$* \\
Llama-3.2-1b vs Qwen2.5-3b-bartowski & 0 & 23 & 0,0 & $<10^{-4}$ & $<10^{-4}$* \\
Llama-3.2-1b vs Qwen2.5-1.5b & 1 & 84 & 0,0 & $<10^{-4}$ & $<10^{-4}$* \\
Llama-3.2-1b vs Qwen2.5-3b & 0 & 24 & 0,0 & $<10^{-4}$ & $<10^{-4}$* \\
Llama-3.2-1b vs Mistral-7B ROMEO (bf16) & 0 & 93 & 0,0 & $<10^{-4}$ & $<10^{-4}$* \\
Mistral-7B vs Qwen2.5-3b-bartowski & 73 & 5 & 14,6 & $<10^{-4}$ & $<10^{-4}$* \\
Mistral-7B vs Qwen2.5-1.5b & 16 & 8 & 2,0 & 0,1516 & 0,7579\phantom{*} \\
Mistral-7B vs Qwen2.5-3b & 73 & 6 & 12,2 & $<10^{-4}$ & $<10^{-4}$* \\
Mistral-7B vs Mistral-7B ROMEO (bf16) & 6 & 8 & 0,8 & 0,7905 & 1,0000\phantom{*} \\
Qwen2.5-3b-bartowski vs Qwen2.5-1.5b & 7 & 67 & 0,1 & $<10^{-4}$ & $<10^{-4}$* \\
Qwen2.5-3b-bartowski vs Qwen2.5-3b & 3 & 4 & 0,8 & 1,0000 & 1,0000\phantom{*} \\
Qwen2.5-3b-bartowski vs Mistral-7B ROMEO (bf16) & 6 & 76 & 0,1 & $<10^{-4}$ & $<10^{-4}$* \\
Qwen2.5-1.5b vs Qwen2.5-3b & 65 & 6 & 10,8 & $<10^{-4}$ & $<10^{-4}$* \\
Qwen2.5-1.5b vs Mistral-7B ROMEO (bf16) & 10 & 20 & 0,5 & 0,0987 & 0,5924\phantom{*} \\
Qwen2.5-3b vs Mistral-7B ROMEO (bf16) & 6 & 75 & 0,1 & $<10^{-4}$ & $<10^{-4}$* \\
\bottomrule
\end{tabular}
\end{table}


\begin{figure}[t]
\centering
\figpng{fig_llm_exactitude}{1}
\caption{Exactitude TrafScript exacte par système (IC Wilson 95\,\%),
globale et par strate. La chaîne évaluée est celle de production
(prompt + KB OACI + validation déterministe).}
\label{fig:llmacc}
\end{figure}

\begin{figure}[t]
\centering
\figpng{fig_llm_categories}{1}
\caption{Exactitude par catégorie de clairance et par système.}
\label{fig:llmcat}
\end{figure}

\begin{figure}[t]
\centering
\figpng{fig_llm_latence}{1}
\caption{Latence d'interprétation et débit de génération des modèles locaux
(RTX 4070, llama.cpp Q4\_K\_M, température 0).}
\label{fig:llmlat}
\end{figure}

\paragraph{Génération de situations.} Sur les 20 descriptions instructeur de
la validation historique (116 contraintes vérifiables : effectifs, niveaux,
caps, azimuts, espacements, types), le générateur à règles atteint
\nScenParser{} et le meilleur LLM local \nScenBest{}
(figure~\ref{fig:llmscen}) : la tâche, plus ouverte (géométrie du secteur,
espacement en traînée), reste sensiblement plus difficile que l'extraction
d'ordres pour les petits modèles.

\begin{figure}[t]
\centering
\figpng{fig_llm_scenarios}{.8}
\caption{Conformité aux contraintes du générateur de situations
(IC Wilson 95\,\%).}
\label{fig:llmscen}
\end{figure}

\subsection{Synthèse vocale (TTS)}
\label{sec:res-tts}

Le tableau~\ref{tab:tts} et la figure~\ref{fig:tts} évaluent Kokoro-82M
(4 voix), Windows SAPI et les voix de pilotes clonées XTTS de la façade
cluster sur les collationnements produits par le générateur de phraséologie
du projet. Kokoro synthétise à RTF \nTTSKokoroRTF{} et son intelligibilité
aller-retour (WER d'un juge STT fixe) reste à \nTTSKokoroWERVHF{}
\emph{après} dégradation VHF : la coloration radio coûte peu
d'intelligibilité, conformément à l'objectif (réalisme sans perte
pédagogique). Les voix clonées XTTS (timbre « pilote » volontairement
réaliste, débit rapide) synthétisent à RTF $\approx$0,78 mais leur WER
aller-retour est plusieurs fois supérieur (tableau~\ref{tab:tts}) : le
réalisme du timbre se paie en intelligibilité mesurée --- compromis
immersion/clarté assumé et \emph{réglable}, le choix du fournisseur TTS
étant une ligne de configuration. La métrique étant relative (biais du juge
partagé), seuls les écarts entre moteurs sont interprétables.

\begin{table}[t]\centering\small
\caption{TTS : RTF et intelligibilité aller-retour (juge faster-whisper small (cuda, greedy)) sur 30 collationnements OACI.}
\label{tab:tts}
\begin{tabular}{lrrr}
\toprule
Moteur / voix & RTF & WER propre & WER après VHF \\
\midrule
kokoro/af\_bella & 0,348 & 22,5\,\% & 25,4\,\% \\
kokoro/am\_adam & 0,346 & 18,2\,\% & 17,7\,\% \\
kokoro/bf\_emma & 0,339 & 21,2\,\% & 19,8\,\% \\
kokoro/bm\_george & 0,310 & 26,0\,\% & 22,6\,\% \\
sapi/Microsoft Zira Desktop & 0,003 & 16,0\,\% & 18,2\,\% \\
\bottomrule
\end{tabular}
\end{table}


\begin{figure}[t]
\centering
\figpng{fig_tts}{1}
\caption{TTS : coût de synthèse et intelligibilité objective aller-retour,
avant/après canal VHF.}
\label{fig:tts}
\end{figure}

\subsection{Boucle vocale complète (E2E)}
\label{sec:res-e2e}

La figure~\ref{fig:e2e} et le tableau~\ref{tab:e2e} mesurent la boucle
complète de production sur \nEEn{} clairances parlées (canal radio SNR
12~dB) ; le tableau~\ref{tab:e2ecmp} compare ensuite cette condition de
référence aux conditions cluster (\S\ref{sec:res-romeo}) et locuteur humain
(\S\ref{sec:res-humain}) sur les mêmes énoncés. En local :
\nEEVoix{} d'exécutions exactes en condition voix, contre
\nEETexte{} en condition témoin texte (mêmes clairances, sans STT) ; l'écart chiffre le coût résiduel de l'étage STT, et l'attribution des échecs
(transcription déviée vs interprétation) est donnée dans le
tableau~\ref{tab:e2e}. La latence vocale totale (STT $+$ LLM $+$ TTS,
dégradation VHF cliente comprise) est de \nEELatMoy{} en moyenne
(p95 \nEELatPceQuinze), compatible avec le rythme d'un exercice
d'entraînement, l'étage LLM dominant le budget.

\begin{table}[t]\centering\small
\caption{Boucle vocale complète : configuration, réussite, attribution des échecs et latences par étage.}
\label{tab:e2e}
\begin{tabular}{lr}
\toprule
Clairances parlées (canal SNR 12\,dB) & 25 \\
Réussite voix (IC Wilson) & 76\,\% \ic{57}{89} \\
Réussite témoin texte & 88\,\% \\
Échecs attribués STT / LLM / fournisseur & 3 / 3 / 0 \\
\midrule
Latence STT (moy. / p95) & 0,36 / 0,42\,s \\
Latence LLM (moy. / p95) & 1,06 / 1,42\,s \\
Latence TTS (moy. / p95) & 1,15 / 1,36\,s \\
Latence totale (moy. / p95) & 2,57 / 3,20\,s \\
\bottomrule
\end{tabular}
\end{table}

\begin{table}[t]\centering\small
\caption{Boucle vocale complète : conditions comparées (25 clairances identiques, graine de bruit commune ; IC de Wilson 95\,\% ; attribution des échecs par WER sémantique ; latences client, tunnel compris pour ROMEO).}
\label{tab:e2ecmp}
\begin{tabular}{lrrrr}
\toprule
Condition & Réussite & IC 95\,\% & Échecs STT/LLM & Lat. moy./p95 \\
\midrule
SAPI $\rightarrow$ pile locale & 19/25 (76\,\%) & \ic{57}{89} & 3/3 & 2,6/3,2\,s \\
SAPI $\rightarrow$ pile ROMEO (tunnel) & 19/25 (76\,\%) & \ic{57}{89} & 3/3 & 9,0/13,6\,s \\
Locuteur humain (nicolas), micro brut $\rightarrow$ pile locale & 20/25 (80\,\%) & \ic{61}{91} & 4/1 & 2,3/2,7\,s \\
Locuteur humain (nicolas), canal radio simulé $\rightarrow$ pile locale & 18/25 (72\,\%) & \ic{52}{86} & 6/1 & 2,7/3,6\,s \\
\bottomrule
\end{tabular}
\end{table}


\begin{figure}[t]
\centering
\includegraphics[width=\linewidth]{figures/fig_e2e}
\caption{Boucle vocale complète : taux de réussite (voix vs témoin texte,
IC Wilson) et décomposition de la latence.}
\label{fig:e2e}
\end{figure}

\subsection{Comparaison cluster/local à protocole identique (ROMEO)}
\label{sec:res-romeo}

À la réouverture du supercalculateur (août 2026), la façade GH200
(Whisper-LoRA, Mistral-7B \emph{bf16 non quantifié}, XTTS voix clonées) a
été soumise, à travers le tunnel SSH de production, aux \emph{mêmes}
protocoles seedés que la pile locale : mêmes 116 clairances, mêmes
échantillons ATCO2/UWB, mêmes 25 énoncés E2E, même chaîne cliente. Les
lignes « ROMEO » des tableaux~\ref{tab:stt}, \ref{tab:llm}, \ref{tab:tts}
et~\ref{tab:e2ecmp} en donnent les résultats ; trois lectures :

\begin{enumerate}
  \item \textbf{Effet de la quantification isolé (RQ3).} Mistral-7B bf16
    (ROMEO) atteint \nLLMRomeoAcc{} contre \nLLMBestAcc{} pour le même
    modèle en Q4\_K\_M local ; le test de McNemar apparié sur la paire
    (\nMcnQuantDetail{}, tableau~\ref{tab:mcnemar}) ne détecte \emph{pas}
    d'écart significatif à cette taille d'échantillon : sur cette tâche, la
    quantification 4~bits est au plus un effet de l'ordre du point de
    pourcentage --- le facteur de confusion « précision numérique » de la
    campagne locale est ainsi levé, dans les limites de puissance
    discutées en \S\ref{sec:validite}.
  \item \textbf{Parité STT.} Le même adaptateur LoRA servi par la façade
    GH200 produit, sur les mêmes extraits, des WER cohérents avec
    l'exécution locale (tableau~\ref{tab:stt}, ligne \syst{romeo-lora}) :
    le point de déploiement ne change pas la qualité de transcription,
    seulement la latence (tunnel compris).
  \item \textbf{Boucle complète.} En bout en bout, la pile ROMEO exécute
    \nEERomeoVoix{} des clairances parlées (témoin texte \nEERomeoTexte)
    pour une latence vocale moyenne de \nEERomeoLatMoy, à comparer à
    \nEEVoix{} / \nEELatMoy{} pour la pile locale
    (tableau~\ref{tab:e2ecmp}). Les latences ROMEO incluent l'aller-retour
    du tunnel SSH : c'est la réalité opérationnelle du poste de travail.
\end{enumerate}

Fait opérationnel documenté : pendant la campagne locale (juillet 2026), le
cluster était \emph{entièrement indisponible} (maintenance de sécurité,
correctifs RedHat) ; l'application est restée fonctionnelle par simple
bascule du \syst{.env} vers la façade locale, puis a rebasculé vers la
façade GH200 pour la présente campagne --- validation en conditions
réelles, dans les deux sens, de l'argument de continuité de service du
contrat OpenAI-compatible.

\subsection{Validation avec locuteur humain}
\label{sec:res-humain}

Les énoncés synthétiques (SAPI) contrôlent parfaitement le protocole mais ne
représentent ni l'accent ni la prosodie d'un utilisateur réel. Un locuteur
humain (francophone, non contrôleur) a donc enregistré les 25 clairances du
protocole E2E via le micro de l'application, rejouées dans la chaîne exacte
du banc en deux conditions : \emph{micro brut} (ce que l'application entend
au push-to-talk) et \emph{canal radio simulé} (bruit SNR 12~dB +
passe-bande, graine commune avec les autres conditions). Résultats au
tableau~\ref{tab:e2ecmp} : \nHumMicro{} de réussite au micro brut
(IC \nHumMicroIC) et \nHumRadio{} après canal radio, contre \nEEVoix{} pour
la voix synthétique --- la boucle tient face à une voix réelle non native,
au prix d'échecs supplémentaires attribués à l'étage STT. L'analyse
qualitative des échecs (\S\ref{sec:discussion}) est plus instructive que
les taux : accent sur certains indicatifs (« ryanair »), confusion
phonétique \emph{two}/\emph{to} transformant « two five zero » en
« to five zero » (soit SPD~50, \emph{dans les bornes} donc non filtré), et
un défaut indépendant du locuteur --- l'indicatif épelé « csa one delta
zulu » parfois réassemblé fautivement par le LLM (\syst{CSAD12Z} au lieu de
\syst{CSA1DZ}) malgré une transcription parfaite, défaut également observé
en condition synthétique.

% =============================================================================
\section{Discussion}
\label{sec:discussion}

\paragraph{Choix du point de déploiement.} Les trois points de déploiement
mesurés dessinent une matrice simple. \emph{Cluster (GH200, bf16)} :
meilleure qualité d'interprétation (\nLLMRomeoAcc) et voix clonées, mais
dépendance à une infrastructure mutualisée dont l'indisponibilité de
juillet 2026 a rappelé le coût, et latences tributaires du tunnel.
\emph{Local (RTX 4070, Q4\_K\_M)} : qualité proche (\nLLMBestAcc{} pour le
même Mistral-7B), latence vocale moyenne de \nEELatMoy, aucune dépendance
réseau, coût matériel unique de l'ordre d'un poste de travail --- le point
recommandé pour un produit d'entraînement autonome. \emph{Cloud} : couvert
par le contrat mais volontairement non évalué (biais commercial,
variabilité temporelle) ; le banc est prêt à l'y soumettre. Dans tous les
cas, les gardes déterministes restent identiques : la qualité du modèle ne
conditionne que le taux d'ordres \emph{corrects}, jamais la sûreté des
ordres \emph{exécutés} (RQ4 : zéro ordre hors bornes exécuté sur l'ensemble
des campagnes, tous modèles confondus).

\paragraph{Taxonomie des échecs résiduels.} Les échecs observés en boucle
complète se rangent en quatre classes. (i)~\emph{Acoustiques} : accents et
indicatifs peu fréquents (« ryanair » humain) --- traitables par
élargissement du corpus d'adaptation LoRA. (ii)~\emph{Phonétiquement
ambigus mais dans les bornes} : « two five zero » entendu « to five zero »
produit SPD~50, valide pour le validateur (bornes 0--350~kt) mais
opérationnellement invraisemblable en croisière --- un plancher de
plausibilité par phase de vol (borne basse dynamique) fermerait cette
classe sans toucher au contrat. (iii)~\emph{Réassemblage d'indicatifs
épelés par le LLM} : \syst{CSAD12Z} au lieu de \syst{CSA1DZ} malgré une
transcription parfaite, observé avec voix humaine \emph{et} synthétique ---
défaut du modèle, pas du locuteur ; une normalisation déterministe des
indicatifs épelés \emph{en amont} du LLM (le module \syst{atc\_callsign}
existe déjà côté post-traitement) le neutraliserait. (iv)~\emph{Style de
sortie du modèle} : l'anomalie Qwen2.5-3B (schéma violé) et la perte du
message système par un gabarit de conversation (Mistral/llama.cpp),
détectées uniquement parce que l'évaluation traverse la pile réelle ---
argument méthodologique central de ce travail.

\paragraph{Portée pour la formation.} Le facteur temps réel du simulateur
($\times$\nScalingDeuxCents{} à 200 aéronefs), la latence vocale
(\nEELatMoy{} locale) et le collationnement dégradé canal radio placent le
système dans l'enveloppe d'un poste d'entraînement individuel réaliste ; le
barème d'exercice, adossé à des conflits garantis par construction
(\nConflitGarantie), fournit une note reproductible et contestable ---
c'est-à-dire auditable, chaque événement étant tracé dans le rapport JSON
de débrief.

% =============================================================================
\section{Menaces à la validité}
\label{sec:validite}

\textbf{Interne.} Les énoncés E2E synthétiques contrôlent le protocole mais
sous-estiment la variabilité de la parole spontanée ; la condition locuteur
humain (\S\ref{sec:res-humain}) borne ce risque, mais avec $n=1$ locuteur
et $n=25$ énoncés --- les IC de Wilson, larges, en rendent compte, et le
recrutement de locuteurs supplémentaires est la première extension prévue.
L'intelligibilité TTS partage le biais du juge STT entre moteurs (métrique
relative). L'erreur CPA mesurée est bornée par l'arrondi d'affichage du
champ comparé. Les latences ROMEO incluent le tunnel SSH (réalité du poste,
mais non séparée de l'inférence). \textbf{Externe.} La comparaison
bf16/Q4\_K\_M isole désormais l'effet de quantification
(\S\ref{sec:res-romeo}) ; les fournisseurs cloud restent non évalués (le
contrat les rend interchangeables, pas équivalents). Le secteur simulé est
unique (CTR Reims, 7 fixes) ; la généralisation à d'autres espaces exige de
régénérer graphe et exercices. \textbf{Construction.} Le corpus
d'interprétation (116 cas) et sa vérité terrain ont été rédigés par un
unique annotateur ; deux garde-fous l'encadrent --- la moitié
« phraséologie standard » est vérifiée mécaniquement par le parseur à
règles (100\,\% d'accord), et l'intégralité du corpus annoté est publiée
avec le code pour contre-expertise. La vérité terrain des cas mixtes
suppose la politique « filtrer l'invalide, garder le valide », qui est
celle du système ; un système plus conservateur serait pénalisé par cette
métrique --- le taux de rejet des négatifs est donné séparément pour cette
raison. \textbf{Puissance.} Avec 116 cas appariés et des effectifs
discordants faibles (typiquement 10--30 paires), le test de McNemar ne
détecte que des écarts de plusieurs points de pourcentage : les paires non
significatives du tableau~\ref{tab:mcnemar} (parseur vs Mistral,
bf16 vs Q4\_K\_M) se lisent « écart au plus modeste », et non
« équivalence démontrée » --- une démonstration d'équivalence exigerait un
protocole dédié (TOST) sur corpus élargi.

% =============================================================================
\section{Travaux liés}

Les corpus ATCOSIM, UWB-ATCC et le challenge ATCO2~\cite{atcosim,uwbatcc,
atco2} ont établi la reconnaissance vocale ATC comme domaine propre, avec
des WER de systèmes génériques dégradés de plusieurs dizaines de points par
le canal VHF ; le retour d'expérience à grande échelle
d'ATCO2~\cite{atco2lessons} confirme le rôle décisif de l'adaptation au
domaine, aujourd'hui dominée par l'affinage paramétriquement efficace
(LoRA~\cite{lora}) de modèles de fondation comme Whisper~\cite{whisper}.
BlueSky~\cite{bluesky} est le simulateur ouvert de référence pour la
recherche ATM. L'utilisation de LLM contraints pour l'interprétation de
clairances rejoint les travaux sur l'extraction structurée fiable ; notre
apport est triple : l'évaluation \emph{système} (à travers la chaîne de
production complète, gardes comprises, jusqu'aux gabarits de conversation
des moteurs d'inférence) plutôt que l'évaluation isolée du modèle ; la
comparaison appariée à un parseur à règles sur corpus stratifié avec
contrôle des comparaisons multiples ; et la triangulation entre points de
déploiement (bf16 cluster / Q4 local) et types de locuteurs (synthétique /
humain) à protocole constant. Les moteurs d'inférence locaux utilisés sont
llama.cpp, CTranslate2/faster-whisper et Kokoro-82M (ONNX), cités ici comme
logiciels\footnote{\syst{github.com/ggml-org/llama.cpp},
\syst{github.com/SYSTRAN/faster-whisper},
\syst{huggingface.co/hexgrad/Kokoro-82M}.}.

% =============================================================================
\section{Conclusion}

Nous avons présenté un simulateur d'entraînement ATC dont chaque étage
(géométrie de conflit, reconnaissance vocale, interprétation, synthèse,
boucle complète) est validé par des mesures contrôlées, seedées et
versionnées avec le code, et dont les réponses aux cinq questions de
recherche sont chiffrées. RQ1 : base géométrique exacte ($10^6$ décisions
sans désaccord, $F_1=1{,}0$ face à BlueSky, conflits d'exercice garantis à
\nConflitGarantie). RQ2 : gain STT majeur et significatif de l'adaptation
au domaine sur VHF réel (\nWERVanillaAtcoVhf{} $\rightarrow$
\nWERLoraAtcoVhf{} sur ATCO2), passe-bande d'inférence validé. RQ3 :
hiérarchie nette des LLM locaux, Mistral-7B Q4 statistiquement
indiscernable du parseur à règles globalement et supérieur hors grammaire ;
la paire bf16/Q4 à modèle constant (\nLLMRomeoAcc{} vs \nLLMBestAcc) borne
le coût de la quantification. RQ4 : zéro ordre hors bornes exécuté sur
l'ensemble des campagnes, quel que soit le modèle. RQ5 : boucle vocale à
\nEEVoix{} de réussite (\nEELatMoy{} de latence moyenne) en local,
\nEERomeoVoix{} via la façade cluster, et \nHumMicro{} avec locuteur humain
réel --- avec une taxonomie d'échecs actionnable. L'architecture au contrat
OpenAI-compatible rend ces points de déploiement interchangeables sans
toucher au code ni aux gardes de sûreté, qui restent déterministes et
testées ; la bascule cluster$\leftrightarrow$local a été exercée dans les
deux sens en conditions réelles.

\paragraph{Reproductibilité.} Code, corpus annotés, scripts de benchmark,
résultats JSON (locaux \emph{et} ROMEO : \syst{bench/results/*\_romeo.json},
locuteur humain : \syst{human\_e2e.json}), environnement logiciel capturé
(tableau~\ref{tab:versions}) et génération automatique des figures et
tableaux de cet article : répertoires \syst{bench/}, \syst{validation/} et
\syst{docs/article/} du dépôt. La campagne historique ré-exécutée produit
des métriques identiques champ à champ.

% =============================================================================
\begin{thebibliography}{9}\small
\bibitem{bluesky} J.~M. Hoekstra et J.~Ellerbroek,
« BlueSky ATC Simulator Project: an Open Data and Open Source Approach »,
\emph{ICRAT}, 2016.
\bibitem{whisper} A.~Radford \emph{et al.},
« Robust Speech Recognition via Large-Scale Weak Supervision »,
\emph{ICML}, 2023.
\bibitem{lora} E.~Hu \emph{et al.}, « LoRA: Low-Rank Adaptation of Large
Language Models », \emph{ICLR}, 2022.
\bibitem{atco2} J.~Zuluaga-Gomez \emph{et al.}, « ATCO2 corpus: A Large-Scale
Dataset for Research on Automatic Speech Recognition and Natural Language
Understanding of Air Traffic Control Communications », 2022.
\bibitem{uwbatcc} L.~Šmídl \emph{et al.}, « Air Traffic Control Communication
(UWB-ATCC) Corpus », Université de Bohême de l'Ouest, 2019.
\bibitem{atcosim} K.~Hofbauer, S.~Petrik et H.~Hering, « The ATCOSIM Corpus of
Non-Prompted Clean Air Traffic Control Speech », \emph{LREC}, 2008.
\bibitem{mistral} A.~Q. Jiang \emph{et al.}, « Mistral 7B », 2023.
\bibitem{qwen} Équipe Qwen, « Qwen2.5 Technical Report », 2024.
\bibitem{atco2lessons} J.~Zuluaga-Gomez \emph{et al.}, « Lessons Learned in
Transcribing 5000~h of Air Traffic Control Communications for Robust
Automatic Speech Understanding », \emph{Aerospace}, 10(10):898, 2023.
\bibitem{mcnemar} Q.~McNemar, « Note on the sampling error of the difference
between correlated proportions or percentages », \emph{Psychometrika},
12(2):153--157, 1947.
\bibitem{holm} S.~Holm, « A simple sequentially rejective multiple test
procedure », \emph{Scandinavian Journal of Statistics}, 6(2):65--70, 1979.
\bibitem{wilson} E.~B. Wilson, « Probable inference, the law of succession,
and statistical inference », \emph{Journal of the American Statistical
Association}, 22(158):209--212, 1927.
\bibitem{efron} B.~Efron, « Bootstrap methods: another look at the
jackknife », \emph{The Annals of Statistics}, 7(1):1--26, 1979.
\bibitem{icao4444} OACI, \emph{Doc 4444 : Gestion du trafic aérien (PANS-ATM)},
16\textsuperscript{e} éd. (phraséologie de référence ; la base de
connaissances du projet est un contenu original, sans extrait verbatim).
\end{thebibliography}

\end{document}
